\section{Cómo trabajan los agentes: macro actions y paralelización}

\subsection{El modelo de Peter Steinberg}

Karpathy mencionó en particular la manera de trabajar de Peter Steinberg: ejecuta simultáneamente varias sesiones de agente de Codex, cada una realizando una tarea de alta dificultad de unos 20 minutos. Alterna entre varios repos y le asigna a cada agente una tarea funcional independiente.

\begin{importantbox}{El pensamiento de macro actions}
El cambio de pensamiento clave es este: la unidad de operación deja de ser «una línea de código» o «una función», y pasa a ser un módulo funcional completo o una línea de investigación, como unidad de macro operación. El Agente 1 hace la función A, el Agente 2 hace la función B (sin que ambas entren en conflicto), el Agente 3 hace investigación y el Agente 4 escribe la propuesta de implementación: todo se gestiona en el repositorio de código mediante macro actions.
\end{importantbox}

\begin{figure}[H]
\centering
\includegraphics[width=0.82\textwidth]{frame-macro-actions.jpg}
\caption{Video aproximadamente 00:04:26: Karpathy explica que, en la era de los code agents, la unidad de gestión del repositorio de software pasó de las modificaciones microscópicas a las macro actions}
\end{figure}

\subsection{Skill issue: todo es un problema de habilidad}

Cuando un agente falla, Karpathy tiende a pensar que se trata de un «skill issue» y no de una falta de capacidad: que él mismo no dio instrucciones suficientemente buenas, no configuró bien el archivo agents.md, o no construyó las herramientas de memory adecuadas. Esta mentalidad de que «todo se puede optimizar» resulta a la vez emocionante y angustiante.

\begin{warningbox}{La jaggedness (lo desparejo) del agente}
Karpathy describió también el lado frustrante de los agentes: siente que está hablando al mismo tiempo con «un doctorado extremadamente inteligente + un niño de 10 años». En los dominios on-rails (como las tareas de programación entrenadas mediante RL) el agente muestra una capacidad sobrehumana, pero en los dominios off-rails (como entender la intención del usuario o contar chistes) se estanca. Esta jaggedness es una característica fundamental de los modelos actuales.
\end{warningbox}

